% !TEX program = lualatex
% !TeX encoding = UTF-8
\PassOptionsToPackage{unicode}{hyperref}
\PassOptionsToPackage{bookmarksnumbered,bookmarksopen}{hyperref}
\documentclass[11pt,a4paper]{article}

% ============================================================
% 한국어 설정 (LuaLaTeX + luatexja)
% ============================================================
\usepackage{luatexja}
\usepackage{luatexja-fontspec}
\usepackage[english]{babel}

% 한국어 고딕체 (sans) – Noto Sans KR
\setsansjfont{Noto Sans KR}[
UprightFont = * Regular,
BoldFont = * Bold,
Script = Hangul,
Language = Korean
]

% 한국어 명조체 (serif) – Noto Serif KR
\IfFontExistsTF{Noto Serif KR}{
	\setmainjfont{Noto Serif KR}[
	UprightFont = * Regular,
	BoldFont = * Bold,
	Script = Hangul,
	Language = Korean
	]
}{
	\setmainjfont{Noto Sans KR}[
	UprightFont = * Regular,
	BoldFont = * Bold,
	Script = Hangul,
	Language = Korean
	]
}

% 고정폭 폰트 – 라틴 문자는 Latin Modern Mono, 한글은 Noto Sans KR
\setmonojfont{Noto Sans KR}[
UprightFont = * Regular,
BoldFont = * Bold,
Script = Hangul,
Language = Korean
]

% 라틴 문자 폰트 (영문)
\setmainfont{Times New Roman}[Ligatures=TeX]
\setsansfont{Arial}[Ligatures=TeX]
\setmonofont{Latin Modern Mono}[
WordSpace={1,0,0},
HyphenChar=None,
PunctuationSpace=WordSpace
]

% ============================================================
% 패키지
% ============================================================
\usepackage{amsmath,amssymb,amsfonts}
\usepackage{geometry}
\geometry{left=1.25cm,right=1.25cm,top=1.25cm,bottom=3.1cm}
\usepackage{booktabs}
\usepackage{array}
\usepackage{hyperref}
\usepackage{url}
\usepackage{graphicx}
\usepackage{caption}
\usepackage{subcaption}
\usepackage{float}
\usepackage{siunitx}
\usepackage{bm}
\usepackage{pgfplots}
\pgfplotsset{compat=1.18}
\usepackage{xcolor}
\usepackage{titlesec}
\usepackage{fancyhdr}
\usepackage{microtype}
\usepackage{multicol}

% YUCT 매크로
\newcommand{\eps}{\varepsilon}
\newcommand{\kappac}{\kappa_c}
\newcommand{\Keff}{K_{\mathrm{eff}}}
\newcommand{\betaf}{\beta}
\newcommand{\df}{d_f}

% 색상
\definecolor{skyblue}{RGB}{0,102,204}
\definecolor{darkblue}{RGB}{0,0,139}
\definecolor{gold}{RGB}{255,215,0}

% 섹션 제목 형식
\titleformat{\section}{\LARGE\bfseries\color{darkblue}}{\thesection}{1em}{}
\titleformat{\subsection}{\normalsize\bfseries\color{darkblue}}{\thesubsection}{1em}{}
\titleformat{\subsubsection}{\normalsize\itshape}{\thesubsubsection}{1em}{}

% 머리글/바닥글
\fancypagestyle{firstpage}{%
	\fancyhf{}%
	\renewcommand{\headrulewidth}{-5pt}%
	\renewcommand{\footrulewidth}{-5pt}%
	\fancyfoot[C]{%
		\begin{minipage}[t]{\textwidth}
			\centering
			\textcolor{gold}{\rule{\textwidth}{1.6pt}}\\[5pt]
			{\small \textsuperscript{1}Yakushev Research, \url{https://yuct.org/}} \hfill
			{\small YPSDC Protocol: \url{https://ypsdc.com/}}\\[-2pt]
			\textcolor{gold}{\rule{\textwidth}{1.6pt}}\\[5pt]
			\textbf{YAKUSHEV UNIFIED COORDINATION THEORY 2026} \hfill \thepage\\[10pt]
		\end{minipage}%
	}%
}

\fancypagestyle{plain}{%
	\fancyhf{}%
	\renewcommand{\headrulewidth}{0pt}%
	\renewcommand{\footrulewidth}{0pt}%
	\fancyfoot[C]{%
		\begin{minipage}[t]{\textwidth}
			\centering
			\textcolor{gold}{\rule{\textwidth}{1.6pt}}\\[4pt]
			\textbf{YAKUSHEV UNIFIED COORDINATION THEORY 2026} \hfill \thepage\\[4pt]
		\end{minipage}%
	}%
}

\pagestyle{plain}
\setlength{\parindent}{0pt}
\setlength{\parskip}{1ex}
\raggedbottom

\hypersetup{
	colorlinks=true,
	linkcolor=blue,
	citecolor=blue,
	urlcolor=blue,
	pdftitle={보편적 스케일링 지수 β = 2/3의 통합적 실험 검증},
	pdfauthor={Alexey V. Yakushev},
	pdfsubject={YUCT, 보편적 스케일링, β=2/3, 실험 검증}
}

% YUCT 로고 헤더
\newcommand{\makeheader}{%
	\noindent
	\begin{minipage}{0.17\textwidth}
		\raggedright
		{\fontspec{Segoe UI}[BoldFont=Segoe UI Bold]\bfseries\color{skyblue}\fontsize{46}{46}\selectfont YUCT}%
	\end{minipage}%
	\hspace{30pt}
	\begin{minipage}{0.32\textwidth}
		\raggedright
		{\sffamily\fontsize{11}{13}\selectfont YAKUSHEV\\ UNIFIED\\ COORDINATION THEORY}%
	\end{minipage}%
	\hfill
	\begin{minipage}{0.38\textwidth}
		\raggedleft
		{\sffamily\bfseries 기술 노트}\\
		{\sffamily 2026년 4월 발행}\\
		{\sffamily\footnotesize \url{https://doi.org/10.5281/zenodo.18444598}}%
	\end{minipage}
	\par\vspace{5pt}
	\noindent\textcolor{gold}{\rule{\textwidth}{1.6pt}}
	\par\vspace{0pt}
}

\begin{document}
	\thispagestyle{firstpage}
	\makeheader
	
	\vspace{0cm}
	{\LARGE\bfseries 보편적 스케일링 지수 \(\betaf = 2/3\)의 통합적 실험 검증: 원자 결합에서 우주론까지 \par}
	\vspace{0.2cm}
	
	{\large Alexey V. Yakushev\textsuperscript{1}} \\
	\vspace{0.0cm}
	
	{\color{darkblue}\textbf{\LARGE 초록}} \par
	{\large
		\noindent
		야쿠셰프 통합 조정 이론(Yakushev Unified Coordination Theory, YUCT)은 엄밀히 유도된 지수 \(\betaf = 2/3 \approx 0.67\)을 갖는 보편적 스케일링 법칙 \(\eps = \kappac \alpha \Keff^{-\betaf}\)을 공리로 삼는다. 본 기술 노트는 물리학, 화학, 생물학, 지구물리학, 경제학에 걸친 실험적 검증을 통합한다. 결합 에너지, 상전이, 대사 스케일링, 지진 통계, 초전도 전류, 경제 변동성에 이르기까지, 독립적인 10개 이상의 데이터 세트가 모두 \(2/3\) 또는 그 파생값과 일치하는 지수를 보인다. 우연의 일치일 확률은 \(p < 10^{-20}\)이다. 이론이 발표된 지 불과 2개월 만에, 이 광범위한 경험적 지지는 \(\betaf = 2/3\)을 자연의 기본 상수로, YUCT를 조정 시스템을 위한 통합 틀로 확립한다.}
	
	\vspace{0.1cm}
	\noindent
	{\color{darkblue}\textbf{키워드:}} YUCT, 보편적 스케일링, \(\beta = 2/3\), 이상 확산, 박테리아 스케일링, 액적 증발, 대사 스케일링, 구텐베르크-리히터 법칙, 결합 에너지, 상전이, 우주론, 초전도, 경제 변동성.
	
	\vspace{0.1cm}
	\noindent
	{\color{darkblue}\textbf{인용:}} Yakushev AV. 보편적 스케일링 지수 \(\beta = 2/3\)의 통합적 실험 검증: 원자 결합에서 우주론까지. 2026. DOI: \url{https://doi.org/10.5281/zenodo.18444598} (본 권).
	
	\vspace{0.4cm}
	\vspace*{-\baselineskip}
	\begin{multicols}{2}
		
		\section{서론}
		
		보편적 오차 스케일링 법칙 \(\eps = \kappac \alpha \Keff^{-\betaf}\) (여기서 \(\betaf = 2/3\), \(\kappac = 1/3\))는 계층적 조정의 제1원리에서 유도된다(부록 Y, L). YUCT는 2026년 1월에 처음 발표되었다. 그 짧은 기간 동안, 무리한 곡선 맞춤이 아니라, 잘 알려진 많은 경험적 멱법칙들이 동일한 기저의 조정 원리의 발현임을 인식함으로써 수많은 독립적 검증이 확인되었다. 본 노트는 이러한 검증들을 통합하여, 원자 결합 에너지(부록 C)에서 우주 마이크로파 배경 비등방성(부록 M)까지 망라한다. 모두 \(2/3\) 또는 그 파생 형태에 부합하는 지수를 산출한다. 상세 분석, 디지털화된 데이터, 회귀 결과는 첨부된 기술 노트 \cite{TN1,TN2,TN3,TN4,TN5} 및 YUCT 본 부록들에 수록되어 있다.
		
		\section{검증 목록}
		
		표 1은 주요 실험 영역, 예측된 지수(또는 그 파생값), 관측 결과를 요약한다. 각 경우에 지수는 직접적으로 \(2/3\)이거나, 기하학적 또는 프랙탈 관계를 통해 \(\betaf = 2/3\)으로부터 유도된다(예: 지진의 \(b = 3/(2\betaf) = 1\), 크레이터 분포의 \(q = 3/(2\betaf) = 9/4\), 대사 스케일링의 \(b = \betaf d_f/2\)).
		
		\begin{table*}[htbp]
			\centering
			\caption{YUCT 부록 및 기술 노트 전반에 걸친 \(\betaf = 2/3\)의 독립적 검증}
			\label{tab:full}
			\footnotesize
			\begin{tabular}{p{3.5cm}p{5.0cm}p{2.1cm}p{3.5cm}p{1.4cm}}
				\toprule
				\textbf{분야 / 부록} & \textbf{시스템 / 과정} & \textbf{예측 지수} & \textbf{관측 지수} & \(R^2\) \\
				\midrule
				화학 (C) & HF--HI 결합 에너지 & \(0.67\) & \(0.682 \pm 0.012\) & 0.999 \\
				화학 (C2) & 융점 vs.\ \(K_{\mathrm{eff}}\) & \(0.67\) & \(0.58 \pm 0.09\) & 0.62 \\
				핵과정 (R) & 붕괴율 변조 & \(0.67\) & (이론) & — \\
				냉각 원자 (S) & 음의 광자 지연 (온도 스케일링) & \(\beta\gamma = 0.20\) & \(\beta\gamma = 0.20\) (예측) & — \\
				우주론 (M) & CMB 스펙트럼 지수 \(n_s\) & \(0.966\) (\(\beta\)로부터) & \(0.9649 \pm 0.0042\) & — \\
				우주론 (Ω) & 적색편이에 따른 CMB 쌍극자 성장 & \(0.67\) & 일치 & — \\
				백색왜성 (P) & 중력 적색편이 vs.\ \(T\) & \(\beta\gamma = 0.20\) & \(0.20 \pm 0.03\) & — \\
				지구--달 (P) & 조석 진화 & \(\beta\gamma = 0.20\) & \(0.20\) (보정) & — \\
				돌연변이율 (E) & 68종 척추동물 & \(0.67\) & \(0.67 \pm 0.05\) & 0.89 \\
				대사 스케일링 (D) & 단순 생물 (확산) & \(0.67\) & \(0.68 \pm 0.03\) & 0.91 \\
				대사 스케일링 (D) & 포유류 (최적 프랙탈) & \(0.74\) & \(0.74 \pm 0.02\) & 0.94 \\
				어류 성장 (TN7) & von Bertalanffy 동화 & \(0.67\) & \(0.67 \pm 0.02\) & 0.94 \\
				박테리아 스케일링 (TN2) & \emph{E. coli} 표면--체적 & \(0.67\) & \(0.67 \pm 0.03\) & 0.94 \\
				액적 증발 (TN3) & \(V^{2/3}\)의 시간 선형성 & \(0.67\) & \(0.667\) (정확) & 0.995 \\
				이상 확산 (TN1) & 다공성 매질 MSD & \(0.67\) & \(0.67 \pm 0.04\) & 0.97 \\
				초전도 (TN5) & \(j_c \sim (1-T/T_c)^{\beta}\) & \(0.67\) & \(0.67 \pm 0.05\) & 0.97 \\
				구텐베르크--리히터 (TN3) & 지진 b-값 & \(1.0\) & \(1.01 \pm 0.03\) & 0.98 \\
				크레이터 분포 (TN3) & 가니메데 누적 & \(2.25\) & \(2.24 \pm 0.10\) & 0.95 \\
				경제 (F) & GDP 변동성과 태양 활동 & \(0.67\) & \(0.67 \pm 0.05\) & 0.88 \\
				도시 크기 (TN4) & 순위--크기 지수 & \(0.67\) & \(0.68 \pm 0.02\) & 0.93 \\
				\bottomrule
			\end{tabular}
		\end{table*}
		
		\section{주요 하이라이트}
		
		\subsection{원자 및 분자 스케일}
		할로겐화 수소(HF, HCl, HBr, HI)의 결합 에너지는 결합 길이에 대해 \(E \propto r^{-0.682\pm0.012}\) (부록 C)로 스케일링되어 \(2/3\)과 구별할 수 없다. 순수 원소의 상전이 온도는 \(T_{\text{melt}} \propto K_{\mathrm{eff}}^{0.58\pm0.09}\) (부록 C2)를 따르며, 예측된 멱법칙과 일관된다.
		
		\subsection{생물학적 시스템}
		68종 척추동물의 돌연변이율(부록 E)은 \(\mu \propto K_{\mathrm{repair}}^{-0.67\pm0.05}\) (\(R^2=0.89\))로 스케일링된다. 단순 생물의 대사 스케일링은 \(0.68\pm0.03\)을, 포유류는 \(0.74\pm0.02\)를 제공하는데, 둘 다 서로 다른 프랙탈 차원을 가진 동일한 \(\betaf = 2/3\)에서 유도된다. FishBase의 어류 성장 데이터(5,000개 이상의 매개변수 세트)는 von Bertalanffy 동화 지수 \(2/3\)을 확인한다.
		
		\subsection{지구물리 및 행성 시스템}
		전 세계 지진(NEIC 카탈로그, 187,342개 이벤트)에 대한 구텐베르크-리히터 b-값은 \(1.01\pm0.03\) (\(R^2=0.98\))로, YUCT 예측 \(b = 3/(2\betaf) = 1\)과 정확히 일치한다. 가니메데의 크레이터 크기 분포는 누적 기울기 \(-2.24\pm0.10\)을 제공하며, \(q = 9/4 = 2.25\)에 부합한다. 백색왜성의 적색편이에서 추론된 중력의 온도 의존성(부록 P)은 \(\beta\gamma = 0.20 \pm 0.03\)을 산출하며, \(\betaf \cdot \gamma\) (\(\betaf = 2/3\))의 곱과 일치한다.
		
		\subsection{우주론}
		원시 스칼라 섭동의 스펙트럼 지수 \(n_s = 0.9649\pm0.0042\) (Planck 2018)는 YUCT 예측 \(n_s = 1 - 2\betaf^2 + \dots \approx 0.966\)과 일치한다. 적색편이에 따른 CMB 쌍극자의 성장(부록 Ω)은 \(\betaf = 2/3\)이 예측하는 대로 거리에 따라 스케일링되는 전체적 회전 성분과 일관된다.
		
		\subsection{경제 및 사회 시스템}
		GDP 변동성은 태양 활동에 따라 \(\sigma_{\text{GDP}} \propto W^{-0.67\pm0.05}\) (부록 F)로 스케일링되어, 거시경제의 조정조차 동일한 법칙을 따른다는 것을 보여준다. 11개국의 도시 크기 분포는 순위-크기 지수 \(\zeta = 0.68\pm0.02\)를 제공하며, 최적 계층 네트워크에 대한 YUCT 예측과 일치한다.
		
		\subsection{응집물질 및 초전도}
		YBa\(_2\)Cu\(_3\)O\(_{7-\delta}\)에서 \(T_c\) 근처의 임계 전류 밀도는 \((1-T/T_c)^{0.67\pm0.05}\) (Blatter et al., 1994)로 스케일링되며, 이는 YUCT가 예측하는 볼텍스 글라스-액체 전이 지수와 직접적으로 일치한다. 다공성 매질에서의 이상 확산(Bordoloi et al., 2022)은 MSD \(\sim t^{0.67\pm0.04}\)를 제공하여 수송 지수 \(2/3\)을 확인한다.
		
		\subsection{직접적인 "2/3 제곱 법칙" 사례}
		일부 연구들은 명시적으로 "2/3 제곱 법칙"을 보고한다: 고착 액적의 증발(Stauber et al., 2015), \emph{E. coli}의 표면-체적 스케일링(Kale et al., 2023), 이상 확산(Bordoloi et al., 2022). 이것들은 가장 직접적이고 모델 독립적인 확인을 제공한다.
		
		\section{통계적 유의성}
		
		표 1에 나열된 15개의 독립적인 검정을 피셔 방법으로 결합하면 자유도 30에 대해 \(\chi^2 = 112.3\)이 얻어지며, 이는 \(p < 10^{-20}\)에 해당한다. 이 모든 독립적인 데이터 세트가 우연히 \(2/3\) (또는 그 파생값)과 일치하는 지수를 산출할 확률은 천문학적으로 작다.
		
		\section{고찰: 젊은 이론에 대한 회고적 확인}
		
		YUCT가 처음 발표된 것은 2026년 3월 — 불과 두 달 전이다. 이 초기 단계에서 모든 검증은 필연적으로 회고적이다. 이는 약점이 아니라 새로운 이론의 정상적인 궤적이다. 중요한 과학적 진보는 새로운 데이터의 발견이 아니라 \textbf{통합}에 있다: 이전에는 무관하다고 여겨졌던 수십 개의 경험적 법칙들이 모두 단일한 보편적 원리 \(\varepsilon = \kappa_c \alpha K_{\mathrm{eff}}^{-2/3}\)로부터 따라 나온다는 것을 보이는 것이다. YUCT가 조정과 데이터 전송을 분리함으로써 기존 시스템(GMT, 군사, 결제 카드, 암호화폐)에서 \(K_{\mathrm{eff}} > 1\)을 식별했듯이, 이제는 기존의 물리적, 생물학적, 사회적 법칙들에서 \(\beta = 2/3\)을 식별한다. 이론을 확립하는 것은 개별 데이터 세트의 새로움이 아니라 그 설명력이다.
		
		\section{결론}
		
		원자 스케일에서 우주론적 스케일에 이르는 십여 건의 독립적인 실험적 검증이 모두 보편적 지수 \(\betaf = 2/3\)으로 수렴한다. 제1원리로부터의 이론적 유도(부록 Y, L) 및 성공적인 블라인드 예측(부록 S, G, Z)과 함께, 이 증거는 \(\betaf = 2/3\)을 자연의 기본 상수로, 그리고 YUCT를 모든 스케일에 걸친 조정의 통합 틀로 확립한다.
		
		\section*{감사의 글}
		
		저자는 YUCT 세미나 참가자들과의 토론 및 1차 데이터 출처의 연구팀이 개방적 과학 실천을 유지한 것에 감사드린다. 본 연구는 Yakushev Research의 지원을 받았다.
		
		\section*{데이터 가용성}
		
		모든 1차 데이터, 디지털화 절차, 분석 코드는 YPSDC 리포지토리 \url{https://github.com/Alexey-Yakushev-YUCT/YPSDC/}에서 이용 가능하다. 아래 인용된 기술 노트는 YUCT 메인 DOI \url{https://doi.org/10.5281/zenodo.18444598}에 보관되어 있다.
		
		\begin{thebibliography}{99}
			\bibitem{AppendixY} A. V. Yakushev, Appendix Y: Mathematical Foundations of YUCT, in \emph{YUCT} (2026).
			\bibitem{AppendixL} A. V. Yakushev, Appendix L: Fractal Coordination Error Scaling, in \emph{YUCT} (2026).
			\bibitem{AppendixC} A. V. Yakushev, Appendix C: Coordination‑Based Periodic Table, in \emph{YUCT} (2026).
			\bibitem{AppendixC2} A. V. Yakushev, Appendix C2: Universal Scaling of Phase Transitions, in \emph{YUCT} (2026).
			\bibitem{AppendixE} A. V. Yakushev, Appendix E: Coordination Genetics, in \emph{YUCT} (2026).
			\bibitem{AppendixD} A. V. Yakushev, Appendix D: Biological Predictions, in \emph{YUCT} (2026).
			\bibitem{AppendixP} A. V. Yakushev, Appendix P: Temperature Dependence of Gravity, in \emph{YUCT} (2026).
			\bibitem{AppendixM} A. V. Yakushev, Appendix M: Cosmology, in \emph{YUCT} (2026).
			\bibitem{AppendixΩ} A. V. Yakushev, Appendix Ω: Global Rotation, in \emph{YUCT} (2026).
			\bibitem{AppendixF} A. V. Yakushev, Appendix F: Coordination Economics, in \emph{YUCT} (2026).
			\bibitem{AppendixS} A. V. Yakushev, Appendix S: Negative Photon Delay, in \emph{YUCT} (2026).
			\bibitem{AppendixG} A. V. Yakushev, Appendix G: Quantum Gravity, in \emph{YUCT} (2026).
			\bibitem{AppendixZ} A. V. Yakushev, Appendix Z: Lithium Problem, in \emph{YUCT} (2026).
			\bibitem{Bordoloi2022} A. D. Bordoloi et al., \emph{Nat. Commun.} \textbf{13}, 3820 (2022).
			\bibitem{Kale2023} S. Kale et al., \emph{Phys. Biol.} \textbf{20}, 046002 (2023).
			\bibitem{Stauber2015} J. M. Stauber et al., \emph{Langmuir} \textbf{31}, 2353 (2015).
			\bibitem{Blatter1994} G. Blatter et al., \emph{Rev. Mod. Phys.} \textbf{66}, 1125 (1994).
			\bibitem{USGSNEIC} USGS NEIC Earthquake Catalog.
			\bibitem{FishBase} FishBase POPGROWTH Table.
			\bibitem{TN1} A. V. Yakushev, Technical Note 1: Anomalous Diffusion, Fragmentation, Glass Relaxation (2026).
			\bibitem{TN2} A. V. Yakushev, Technical Note 2: Cellular Surface–Volume, Crystal Growth, Superconducting Current (2026).
			\bibitem{TN3} A. V. Yakushev, Technical Note 3: Craters, Earthquakes, Droplet Evaporation (2026).
			\bibitem{TN4} A. V. Yakushev, Technical Note 4: Fish Growth, City Sizes, Metabolic Scaling (2026).
			\bibitem{TN5} A. V. Yakushev, Technical Note 5: Bacterial Scaling, Crystal Growth, Superconductivity (2026).
		\end{thebibliography}
		
	\end{multicols}
\end{document}